\documentclass[brazil, 12pt]{article}

\usepackage[portuguese]{babel}
\usepackage[utf8]{inputenc}
\usepackage[T1]{fontenc}
\usepackage[scale=0.8]{geometry}
\usepackage{amsmath}
\usepackage{amssymb}
\usepackage{float}
\usepackage{xcolor}
\usepackage{url}
\usepackage{tikz}
\usetikzlibrary{arrows.meta, calc, positioning}

% Cores usadas na figura (para lembrar a caneta azul/vermelha da anotação)
\definecolor{azul}{RGB}{0,80,200}
\definecolor{vermelho}{RGB}{200,30,30}

% Macro para deixar o lugar das figuras indicado sem quebrar a compilação
% caso a imagem ainda não esteja na pasta do Overleaf.
\newcommand{\figplaceholder}[3]{%
\begin{figure}[H]
    \centering
    \IfFileExists{#1}{%
        \includegraphics[width=0.75\linewidth]{#1}%
    }{%
        \fbox{%
        \begin{minipage}[c][4.5cm][c]{0.75\linewidth}
        \centering
        \textbf{Inserir figura aqui}\par
        \vspace{0.3cm}
        Arquivo sugerido: \texttt{#1}
        \end{minipage}}
    }
    \caption{#2}
    \label{#3}
\end{figure}
}

\begin{document}
\shorthandoff{"}

%-----------------------------------------------------------------------------------------------
%       CABEÇALHO
%-----------------------------------------------------------------------------------------------
\begin{center}
\textbf{Instituto Tecnológico de Aeronáutica - ITA} \\
\textbf{Sistemas de Controle Contínuos e Discretos - CMC-12} \\
\textbf{Alunos}: Pietro Maragno Trindade Marcucci, Bernardo Affonso Cheron Vendramini, Matheus Felipe Ramos Borges
\end{center}

\begin{center}
\textbf{Estudo Comparativo de Estimadores de Estado num Pêndulo Invertido sob
Malha de Controle}
\end{center}
%-----------------------------------------------------------------------------------------------
\vspace*{0.5cm}

\section{Introdução}

Neste trabalho implementamos e comparamos, dentro de uma malha de controle
fechada, quatro técnicas de estimação de estados aplicadas a um pêndulo
invertido sobre carro (\emph{cart-pole}) --- um sistema não-linear e
naturalmente instável.

Em CMC-12 vimos apenas a filtragem passa-baixas clássica, sem estimação por
espaço de estados. Por isso confrontamos quatro níveis de estimadores, do que
foi visto em aula ao estado da arte:

\[
\underbrace{\text{Diferenciação} + \text{passa-baixas}}_{\text{\emph{baseline} do curso}}
\;\rightarrow\;
\underbrace{\text{KF linear}}_{\text{linearizado}}
\;\rightarrow\;
\underbrace{\text{EKF}}_{\text{Jacobiano \emph{online}}}
\;\rightarrow\;
\underbrace{\text{UKF}}_{\text{\emph{sigma points}}}.
\]

Escolhemos o pêndulo invertido porque o controle opera em duas fases que
exercitam os filtros em regimes bem diferentes: o \emph{swing-up} (baseado em
energia) leva a haste por ângulos grandes, num regime fortemente não-linear em
que o KF linear degrada e EKF/UKF se destacam; já a estabilização (LQR na
vertical) opera num regime quase-linear, em que todos os filtros convergem de
forma parecida. Assim, um único experimento mostra onde cada filtro compensa.

O estado é
$X = [\,x\ \ \dot{x}\ \ \theta\ \ \dot{\theta}\,]^{\mathsf T}$ (posição e
velocidade do carro; ângulo e velocidade angular da haste). Só medimos as
posições $x$ e $\theta$, com ruído (encoders); as velocidades $\dot{x}$ e
$\dot{\theta}$ são estimadas. Avaliamos o desempenho por RMSE por estado, tempo
de convergência, robustez (ruído alto, baixa taxa de amostragem, erro de
sintonia de $Q/R$, má inicialização), consistência estatística (NEES/NIS) e
custo computacional, comparando também a malha fechada com estado estimado
\emph{vs.} estado ideal (princípio da separação na prática).

\subsection{Parâmetros adotados}

Centralizamos todas as constantes físicas e de simulação em \texttt{params.m}
e as resumimos na Tabela~\ref{tab:params}. Modelamos a haste como barra
uniforme, de modo que o momento de inércia em torno do centro de massa é
$I = \tfrac{1}{3}mL^{2}$.

\begin{table}[H]
\centering
\begin{tabular}{l l l}
\hline
\textbf{Grandeza} & \textbf{Símbolo} & \textbf{Valor} \\
\hline
Massa do carro          & $M$            & $1{,}0$~kg \\
Massa da haste          & $m$            & $0{,}3$~kg \\
Distância pivô--CM       & $L$            & $0{,}5$~m \\
Inércia da haste (CM)   & $I=\tfrac{1}{3}mL^{2}$ & $0{,}025$~kg$\cdot$m$^{2}$ \\
Atrito viscoso do carro & $b$            & $0{,}1$~N$\cdot$s/m \\
Gravidade               & $g$            & $9{,}81$~m/s$^{2}$ \\
\hline
Passo de amostragem     & $\Delta t$     & $0{,}01$~s ($100$~Hz) \\
Horizonte de simulação  & $T_f$          & $10$~s (padrão) \\
\hline
Ruído do encoder ($x$)  & $\sigma_x$     & $0{,}005$~m \\
Ruído do encoder ($\theta$) & $\sigma_\theta$ & $0{,}005$~rad \\
\hline
\end{tabular}
\caption{Parâmetros físicos, de simulação e de medição adotados no projeto.}
\label{tab:params}
\end{table}

\section{Planta de controle}

Deduzimos as equações de movimento pelo procedimento de Lagrange/Newton para o
cart-pole, seguindo MathWorks~\cite{mathworks} e Tedrake~\cite{mit}.

%-----------------------------------------------------------------------------------------------
%   DIAGRAMA DE CORPO LIVRE NO REFERENCIAL NAO INERCIAL DO CARRINHO
%-----------------------------------------------------------------------------------------------
\noindent{No referencial não inercial do carrinho:}

\vspace{0.4cm}

\begin{center}
\begin{tikzpicture}[
    scale=1.15,
    >=Latex,
    every node/.style={font=\large},
    force/.style={->, thick, line width=1.0pt},
    body/.style={thick, line width=1.1pt},
    rod/.style={line width=3.5pt},
    axis/.style={->, thick, line width=0.9pt},
    measure/.style={<->, thick, line width=0.8pt}
]

    % --- Parametros geometricos -----------------------------------------
    \def\rodang{108}      % angulo da haste medido do eixo x (>90 => inclina para a esquerda)
    \def\rodlen{3.6}      % comprimento desenhado da haste
    \def\comfrac{0.55}    % fracao ate o centro de massa

    % --- Chao -----------------------------------------------------------
    \draw[thick] (-4.6,0) -- (4.8,0);
    \foreach \xx in {-4.4,-3.9,...,4.6}{
        \draw[thin] (\xx,0) -- ++(-0.22,-0.22);
    }

    % --- Carrinho -------------------------------------------------------
    \draw[body, fill=gray!10] (-1.7,0.55) rectangle (1.7,1.45);
    \draw[body, fill=white] (-0.95,0.30) circle (0.24);
    \draw[body, fill=white] (0.95,0.30) circle (0.24);

    % --- Pivo e haste ---------------------------------------------------
    \coordinate (P) at (-0.55,1.45);
    \coordinate (Top) at ($(P)+({\rodlen*cos(\rodang)},{\rodlen*sin(\rodang)})$);
    \coordinate (G)   at ($(P)+({\comfrac*\rodlen*cos(\rodang)},{\comfrac*\rodlen*sin(\rodang)})$);
    \coordinate (V)   at ($(P)+(0,2.9)$);

    \draw[dashed, thick] (P) -- (V);          % vertical de referencia
    \draw[rod] (P) -- (Top);                  % haste
    \fill (P) circle (3.2pt);                 % pivo
    \fill (G) circle (4pt);                   % centro de massa

    % --- Aceleracao angular (theta ponto ponto) -------------------------
    \draw[->, thick] ($(Top)+(0.35,0.05)$) arc[start angle=30, end angle=150, radius=0.55];
    \node at ($(Top)+(0.0,0.65)$) {$\ddot{\theta}$};

    % --- Pseudo-forca m xddot no CM (aponta para tras, esquerda) --------
    \draw[force] (G) -- ++(-1.25,0);
    \node[anchor=south east] at ($(G)+(-1.25,0.02)$) {$m\ddot{x}$};

    % --- Peso mg ---------------------------------------------------------
    \draw[force] (G) -- ++(0,-1.0);
    \node[anchor=north east] at ($(G)+(-0.05,-0.55)$) {$mg$};

    % --- Reacao vertical no pivo Fy (seta um pouco a direita da tracejada)
    \draw[force, azul] ($(P)+(0.10,0)$) -- ++(0,1.30);
    \node[anchor=west, azul] at ($(P)+(0.22,1.28)$) {$F_y$};

    % --- Comprimento L (rotulo a direita da haste, parte baixa) ---------
    \node[anchor=west] at ($(P)!0.35!(G)+(0.30,0)$) {$L$};

    % --- Angulo theta (entre vertical e haste, junto ao pivo) -----------
    \draw[->, thick] ($(P)+(0,1.20)$) arc[start angle=90, end angle=\rodang, radius=1.20];
    \node at ($(P)+(-0.16,1.54)$) {$\theta$};

    % --- Reacao horizontal no pivo Fx (rotulo abaixo da seta, no carrinho)
    \draw[force, azul] (P) -- ++(-1.05,0);
    \node[anchor=north, azul] at ($(P)+(-0.62,-0.04)$) {$F_x$};

    % --- Forca de controle u (no carrinho, aponta para a direita) -------
    \draw[force, vermelho] (-2.90,0.95) -- (-1.7,0.95);
    \node[vermelho] at (-2.35,1.24) {$u$};

    % --- Atrito viscoso b xdot (embaixo, aponta para a esquerda) --------
    \draw[force, vermelho] (-1.4,0.15) -- ++(-1.1,0);
    \node[anchor=south, vermelho] at (-1.95,0.20) {$b\dot{x}$};

    % --- Termo inercial M xddot (no carrinho, aponta para a esquerda) ---
    \draw[force] (2.75,1.05) -- (1.75,1.05);
    \node[anchor=west] at (2.30,1.35) {$M\ddot{x}$};

    % --- Reacao no carrinho Fx (aponta para a direita) ------------------
    \draw[force, vermelho] (1.75,0.70) -- (2.90,0.70);
    \node[anchor=west, vermelho] at (2.92,0.70) {$F_x$};

\end{tikzpicture}
\end{center}

\vspace{0.4cm}

%-----------------------------------------------------------------------------------------------
%   EQUACOES
%-----------------------------------------------------------------------------------------------

\begin{align}
    u + F_x - b\dot{x} &= M\ddot{x} \\[4pt]
    mgL\sin\theta + m\ddot{x}\,L\cos\theta &= \left(I + mL^{2}\right)\ddot{\theta} \\[4pt]
    F_x + m\ddot{x} &= m\left(L\ddot{\theta}\cos\theta - L\dot{\theta}^{2}\sin\theta\right)
\end{align}

Substituindo $F_x$ da equação (3) na equação (1):

\[
u + m\left(-\ddot{x} + L\ddot{\theta}\cos\theta - L\dot{\theta}^{2}\sin\theta\right) - b\dot{x} = M\ddot{x}
\]

Reorganizando, chegamos ao sistema:

\[
\left\{
\begin{aligned}
u &= (M+m)\ddot{x} + b\dot{x} - mL\ddot{\theta}\cos\theta + mL\dot{\theta}^{2}\sin\theta \\[6pt]
mL\ddot{x}\cos\theta &= \left(I + mL^{2}\right)\ddot{\theta} - mgL\sin\theta
\end{aligned}
\right.
\]

\vspace{0.3cm}

%-----------------------------------------------------------------------------------------------
%   DESENVOLVIMENTO ATE A FORMA DE ESTADOS
%-----------------------------------------------------------------------------------------------

Escolhemos como vetor de estado

\[
X =
\begin{bmatrix}
x_1 \\ x_2 \\ x_3 \\ x_4
\end{bmatrix}
=
\begin{bmatrix}
x \\ \dot{x} \\ \theta \\ \dot{\theta}
\end{bmatrix},
\qquad\text{de modo que}\qquad
\dot{X} =
\begin{bmatrix}
\dot{x} \\ \ddot{x} \\ \dot{\theta} \\ \ddot{\theta}
\end{bmatrix}
=
\begin{bmatrix}
x_2 \\ \ddot{x} \\ x_4 \\ \ddot{\theta}
\end{bmatrix}.
\]

Como $\dot{x}_1 = x_2$ e $\dot{x}_3 = x_4$ são imediatos, resta obter expressões
explícitas para $\ddot{x}$ e $\ddot{\theta}$.


Reescrevendo o sistema com os termos de aceleração à esquerda:

\[
\begin{aligned}
(M+m)\,\ddot{x} \;-\; mL\cos\theta\,\ddot{\theta}
&= u - b\dot{x} - mL\dot{\theta}^{2}\sin\theta \\[4pt]
-\,mL\cos\theta\,\ddot{x} \;+\; \left(I+mL^{2}\right)\ddot{\theta}
&= mgL\sin\theta
\end{aligned}
\]

Podemos usar o formato matricial

\[
\underbrace{
\begin{bmatrix}
M+m & -mL\cos\theta \\[2pt]
-mL\cos\theta & I+mL^{2}
\end{bmatrix}}_{A}
\begin{bmatrix}
\ddot{x} \\[2pt] \ddot{\theta}
\end{bmatrix}
=
\begin{bmatrix}
u - b\dot{x} - mL\dot{\theta}^{2}\sin\theta \\[2pt]
mgL\sin\theta
\end{bmatrix}.
\]

Para simplificar, definimos

\[
a = M+m, \qquad
h = mL\cos\theta, \qquad
J = I + mL^{2},
\]
\[
r_1 = u - b\dot{x} - mL\dot{\theta}^{2}\sin\theta, \qquad
r_2 = mgL\sin\theta,
\]

de forma que $A=\begin{bmatrix} a & -h \\ -h & J \end{bmatrix}$.


O determinante é

\[
\Delta = \det A = aJ - h^{2}
= (M+m)\left(I+mL^{2}\right) - m^{2}L^{2}\cos^{2}\theta,
\]

e a inversa de uma matriz $2\times2$ simétrica é

\[
A^{-1} = \frac{1}{\Delta}
\begin{bmatrix}
J & h \\ h & a
\end{bmatrix}.
\]

Logo,

\[
\begin{bmatrix}
\ddot{x} \\[2pt] \ddot{\theta}
\end{bmatrix}
= \frac{1}{\Delta}
\begin{bmatrix}
J & h \\ h & a
\end{bmatrix}
\begin{bmatrix}
r_1 \\[2pt] r_2
\end{bmatrix}
\quad\Longrightarrow\quad
\begin{aligned}
\ddot{x} &= \frac{J\,r_1 + h\,r_2}{\Delta}, \\[6pt]
\ddot{\theta} &= \frac{h\,r_1 + a\,r_2}{\Delta}.
\end{aligned}
\]


Substituindo $a$, $h$, $J$, $r_1$ e $r_2$:

\[
\ddot{x} =
\frac{
\left(I+mL^{2}\right)\!\left(u - b\dot{x} - mL\dot{\theta}^{2}\sin\theta\right)
+ mL\cos\theta\,\big(mgL\sin\theta\big)
}{
(M+m)\left(I+mL^{2}\right) - m^{2}L^{2}\cos^{2}\theta
},
\]

\[
\ddot{\theta} =
\frac{
mL\cos\theta\,\!\left(u - b\dot{x} - mL\dot{\theta}^{2}\sin\theta\right)
+ (M+m)\,\big(mgL\sin\theta\big)
}{
(M+m)\left(I+mL^{2}\right) - m^{2}L^{2}\cos^{2}\theta
}.
\]

Para a haste uniforme, $I=\dfrac{mL^{2}}{3}$, de modo que $I+mL^{2}=\dfrac{4}{3}mL^{2}$.


Reunindo tudo, com $X=[\,x\ \ \dot{x}\ \ \theta\ \ \dot{\theta}\,]^{\mathsf T}$:

\[
\dot{X} =
\begin{bmatrix}
\dot{x}_1 \\[4pt] \dot{x}_2 \\[4pt] \dot{x}_3 \\[4pt] \dot{x}_4
\end{bmatrix}
=
\begin{bmatrix}
x_2 \\[6pt]
\dfrac{
\left(I+mL^{2}\right)\!\left(u - b\,x_2 - mL\,x_4^{2}\sin x_3\right)
+ m^{2}gL^{2}\sin x_3\cos x_3
}{
(M+m)\left(I+mL^{2}\right) - m^{2}L^{2}\cos^{2} x_3
} \\[16pt]
x_4 \\[6pt]
\dfrac{
mL\cos x_3\,\!\left(u - b\,x_2 - mL\,x_4^{2}\sin x_3\right)
+ (M+m)\,mgL\sin x_3
}{
(M+m)\left(I+mL^{2}\right) - m^{2}L^{2}\cos^{2} x_3
}
\end{bmatrix}.
\]

\vspace{0.3cm}

%-----------------------------------------------------------------------------------------------
%   LINEARIZACAO EM TORNO DA VERTICAL
%-----------------------------------------------------------------------------------------------
% >>> [SEÇÃO B] label adicionado só para permitir referência cruzada a partir da
% nova Seção "Estimadores de Estado Clássicos" (\ref{sec:linearizacao} logo abaixo);
% o conteúdo desta seção não foi alterado.
\section{Linearização em torno da vertical}
\label{sec:linearizacao}

O controlador de estabilização deve manter a haste na vertical ($\theta = 0$).
Linearizamos a dinâmica em torno desse equilíbrio

\[
X_{eq} = \begin{bmatrix} 0 & 0 & 0 & 0 \end{bmatrix}^{\mathsf T},
\qquad u_{eq} = 0,
\]

usando as aproximações de pequenos ângulos e desprezando os termos de segunda
ordem (o produto $\dot{\theta}^{2}\sin\theta$ é de ordem superior e se anula):

\[
\cos\theta \approx 1, \qquad \sin\theta \approx \theta, \qquad
\dot{\theta}^{2}\sin\theta \approx 0.
\]

Nessas condições o determinante deixa de depender de $\theta$ (sua dependência é
em $\cos^{2}\theta$, cuja derivada em $\theta=0$ é nula):

\[
\Delta_0 = (M+m)\left(I+mL^{2}\right) - m^{2}L^{2}
= mL^{2}\!\left[\tfrac{4}{3}(M+m) - m\right].
\]

As expressões de $\ddot{x}$ e $\ddot{\theta}$ tornam-se lineares:

\[
\begin{aligned}
\ddot{x} &\approx
-\,\frac{\left(I+mL^{2}\right)b}{\Delta_0}\,\dot{x}
\;+\; \frac{m^{2}L^{2}g}{\Delta_0}\,\theta
\;+\; \frac{I+mL^{2}}{\Delta_0}\,u, \\[8pt]
\ddot{\theta} &\approx
-\,\frac{mLb}{\Delta_0}\,\dot{x}
\;+\; \frac{(M+m)mgL}{\Delta_0}\,\theta
\;+\; \frac{mL}{\Delta_0}\,u.
\end{aligned}
\]

Escrevendo $\dot{X} = A\,X + B\,u$, obtemos a \emph{matriz de estados} $A$ e a
\emph{matriz de entrada} $B$:

\[
A =
\begin{bmatrix}
0 & 1 & 0 & 0 \\[4pt]
0 & -\dfrac{\left(I+mL^{2}\right)b}{\Delta_0} & \dfrac{m^{2}L^{2}g}{\Delta_0} & 0 \\[10pt]
0 & 0 & 0 & 1 \\[4pt]
0 & -\dfrac{mLb}{\Delta_0} & \dfrac{(M+m)mgL}{\Delta_0} & 0
\end{bmatrix},
\qquad
B =
\begin{bmatrix}
0 \\[4pt] \dfrac{I+mL^{2}}{\Delta_0} \\[10pt] 0 \\[4pt] \dfrac{mL}{\Delta_0}
\end{bmatrix}.
\]

O termo $A_{4,3} = (M+m)mgL/\Delta_0 > 0$ confirma que a vertical é um equilíbrio
\textbf{instável}: um desvio $\theta>0$ gera $\ddot{\theta}>0$, afastando ainda
mais a haste. É esse comportamento que a realimentação LQR corrige.

\medskip

\noindent\textbf{Instabilidade pelo critério de Routh-Hurwitz.}
Isolando o modo da haste (equação de $\ddot{\theta}$ linearizada), desprezando o
acoplamento com o carro e o atrito e tomando $u=0$, obtemos
$\ddot{\theta} - A_{4,3}\,\theta = 0$, cujo polinômio característico é

\[
P(s) = s^{2} + 0\cdot s - A_{4,3},
\qquad A_{4,3} = \frac{(M+m)mgL}{\Delta_0} = 17{,}795 > 0.
\]

A condição necessária de Routh-Hurwitz já falha em dois pontos: falta o
coeficiente de $s^{1}$ ($a_1=0$) e o termo independente $a_2=-A_{4,3}$ tem sinal
oposto ao de $a_0=1$. Logo $P(s)$ não pode ter todas as raízes no semiplano
esquerdo.

Montando o arranjo de Routh (com $a_1\to\varepsilon\to0^+$ para o zero na
primeira coluna):

\[
\begin{array}{c|cc}
s^{2} & 1 & -A_{4,3} \\[2pt]
s^{1} & \varepsilon & 0 \\[2pt]
s^{0} & -A_{4,3} &
\end{array}
\]

A primeira coluna é $1,\ \varepsilon,\ -A_{4,3}$, ou seja $(+),(+),(-)$: há uma
troca de sinal, logo uma raiz no semiplano direito. As raízes exatas são
$s=\pm\sqrt{A_{4,3}}=\pm4{,}22$, e é a raiz $+4{,}22$ ($e^{+4{,}22\,t}$) que faz
a haste cair. É esse polo instável que a realimentação LQR reposiciona no
semiplano esquerdo.

\medskip

\noindent\textbf{Valores numéricos.} Com
$M=1{,}0$~kg, $m=0{,}3$~kg, $L=0{,}5$~m, $g=9{,}81$~m/s$^2$, $b=0{,}1$~N$\cdot$s/m
e $I=mL^{2}/3=0{,}025$~kg$\cdot$m$^2$, tem-se $I+mL^{2}=0{,}1$ e
$\Delta_0 = 0{,}1075$, resultando em

\[
A =
\begin{bmatrix}
0 & 1 & 0 & 0 \\
0 & -0{,}0930 & 2{,}0533 & 0 \\
0 & 0 & 0 & 1 \\
0 & -0{,}1395 & 17{,}795 & 0
\end{bmatrix},
\qquad
B =
\begin{bmatrix}
0 \\ 0{,}9302 \\ 0 \\ 1{,}3953
\end{bmatrix}.
\]

%-----------------------------------------------------------------------------------------------
%   PROJETO DO LQR
%-----------------------------------------------------------------------------------------------
\section{Projeto do controlador LQR}

Escolhemos o LQR porque a planta é instável em malha aberta (Seção anterior,
Routh-Hurwitz): sem realimentação a haste cai, então um controlador ativo é
obrigatório. Além disso, o cart-pole tem quatro estados fortemente acoplados e
uma única entrada ($u$), o que torna um PID por variável difícil de sintonizar
--- a mesma força precisa regular carro e haste ao mesmo tempo. A realimentação
de estados $u=-K\,X$ ataca os quatro estados de uma vez, e o LQR dá um critério
sistemático para os ganhos de $K$: minimiza um índice de desempenho que
equilibra erro de regulação ($Q$) e esforço de controle ($R$), em vez de
posicionar polos por tentativa. Desde que $(A,B)$ seja controlável, o LQR
garante malha fechada estável com boas margens de robustez. Por fim, o LQR se
encaixa no tema do trabalho: pelo princípio da separação, o mesmo ganho pode
operar sobre o estado estimado ($u=-K\,\hat X$), permitindo comparar a malha
alimentada por cada estimador (KF, EKF, UKF).

\subsection{Controlabilidade}

Antes de projetar, verificamos que o par $(A,B)$ é controlável, isto é, que a
matriz de controlabilidade tem posto completo:

\[
\mathcal{C} = \begin{bmatrix} B & AB & A^{2}B & A^{3}B \end{bmatrix},
\qquad \operatorname{posto}(\mathcal{C}) = 4.
\]

Para o cart-pole isso se verifica: a força aplicada ao carro é capaz de
governar os quatro estados.

\subsection{Formulação do problema}

O regulador linear quadrático escolhe a lei de realimentação de estados
$u = -K\,X$ que minimiza o funcional de custo

\[
J = \int_{0}^{\infty}
\left( X^{\mathsf T} Q\, X + u^{\mathsf T} R\, u \right) dt,
\qquad Q = Q^{\mathsf T} \succeq 0, \quad R = R^{\mathsf T} \succ 0,
\]

em que $Q$ pondera o erro de estado e $R$ pondera o esforço de controle. O ganho
ótimo é

\[
K = R^{-1} B^{\mathsf T} P,
\]

onde $P = P^{\mathsf T} \succeq 0$ é a solução da equação algébrica de Riccati (ARE)

\[
A^{\mathsf T} P + P A - P B R^{-1} B^{\mathsf T} P + Q = 0.
\]

Para um sistema de quarta ordem não resolvemos a ARE à mão; obtemos $P$ (e daí
$K$) numericamente (Seção~\ref{sec:matlab}).

\subsection{Escolha de $Q$ e $R$ (regra de Bryson)}

Como ponto de partida, aplicamos a regra de Bryson: cada peso é o inverso do
quadrado do desvio máximo aceitável para aquela variável,

\[
Q_{ii} = \frac{1}{x_{i,\max}^{2}}, \qquad R = \frac{1}{u_{\max}^{2}}.
\]

Adotando os limites físicos $x_{\max}=0{,}5$~m, $\dot{x}_{\max}=2$~m/s,
$\theta_{\max}=10^{\circ}\approx0{,}175$~rad, $\dot{\theta}_{\max}=2$~rad/s e
$u_{\max}=20$~N:

\[
Q = \operatorname{diag}(4;\ 0{,}25;\ 32{,}8;\ 0{,}25),
\qquad R = 0{,}0025.
\]

Esses valores são só o ponto de partida; a sintonia final foi iterativa
(aumentar $Q_{33}$ deixa o ângulo mais firme; aumentar $R$ deixa a resposta mais
suave e econômica em força).

\subsection{Cálculo do ganho no MATLAB}
\label{sec:matlab}

Com $A$, $B$, $Q$ e $R$ definidos, obtivemos o ganho pela função \texttt{lqr}
do MATLAB (\emph{Control System Toolbox}), reaproveitando o Jacobiano analítico
já implementado em \texttt{jacobian\_f.m} para montar $A$ e $B$ em $\theta=0$.
Verificamos a controlabilidade de $(A,B)$ antes de resolver a equação de
Riccati; em seguida \texttt{lqr(A, B, Q, R)} retorna o ganho ótimo $K$, e
\texttt{eig(A - B*K)} dá os polos de malha fechada.

O ganho obtido é

\[
K = \begin{bmatrix} -40{,}00 & -42{,}14 & 217{,}20 & 48{,}46 \end{bmatrix},
\]

com polos de malha fechada

\[
\begin{aligned}
s_{1,2} &= -12{,}8126 \pm 2{,}5904\,j, \\
s_{3,4} &= -1{,}4435 \pm 1{,}0585\,j,
\end{aligned}
\]

todos com parte real negativa, o que confirma a estabilização da vertical.

%-----------------------------------------------------------------------------------------------
%   SWING-UP ENERGETICO
%-----------------------------------------------------------------------------------------------
\section{Controle de \emph{swing-up} energético}

O LQR só estabiliza a haste perto da vertical (sua região de atração). Quando
o pêndulo parte pendurado, é preciso levá-lo até essa vizinhança primeiro.
Como a linearização não vale para ângulos grandes, usamos uma estratégia
não-linear baseada em energia (\emph{energy shaping}), seguindo Tedrake~\cite{mit}:
em vez de seguir uma trajetória, injetamos energia no pêndulo balançando o
carro no ritmo certo, até ele ter energia suficiente para alcançar a vertical.

\subsection{Energia do pêndulo}

Adotamos a notação de~\cite{mit}. A energia da haste é a soma das parcelas
cinética $T$ e potencial $U$,

\[
E(\theta,\dot{\theta}) = T(\dot{\theta}) + U(\theta),
\qquad
T = \tfrac{1}{2} I\,\dot{\theta}^{2},
\qquad
U = -\,m g l\cos\theta,
\]

onde $I$ é o momento de inércia efetivo em torno do pivô e $l$ a distância do
pivô ao centro de massa. Para a haste uniforme deste trabalho,
$I = \tfrac{1}{3}mL^{2} + mL^{2} = \tfrac{4}{3}mL^{2}$ e $l = L$.

\noindent\emph{Observação de convenção.} A referência~\cite{mit} mede $\theta$ a
partir da posição pendurada ($\theta=0$ embaixo), com $U=-mgl\cos\theta$ mínimo
embaixo. Aqui $\theta=0$ é a vertical superior; as duas convenções diferem só
por uma defasagem de $\pi$ e um sinal no termo potencial, sem alterar a lei de
controle.

A energia da órbita que alcança o equilíbrio superior instável (energia
desejada) é, como em~\cite{mit},

\[
E^{d} = m g l.
\]

O objetivo do \emph{swing-up} é levar $E$ até $E^{d}$. Define-se o erro de energia

\[
\tilde{E} = E - E^{d}.
\]

\subsection{Lei de \emph{energy shaping}}

Para o pêndulo simples atuado por torque na junta, \cite{mit} mostra que a lei

\[
u = -\,k\,\dot{\theta}\,\tilde{E}, \qquad k>0,
\]

produz a dinâmica de erro $\dot{\tilde{E}} = -\,k\,\dot{\theta}^{2}\,\tilde{E}$,
que faz $\tilde{E}\to 0$. Isso decorre da função de Lyapunov
$V = \tfrac{1}{2}\tilde{E}^{2}\ge 0$, cuja derivada é

\[
\dot{V} = \tilde{E}\,\dot{\tilde{E}}
        = -\,k\,\dot{\theta}^{2}\,\tilde{E}^{2} \le 0.
\]

No cart-pole a entrada não é um torque na junta, mas a força $u$ no carro, que
atua sobre a haste pela aceleração do pivô $\ddot{x}$. Da equação de movimento
da haste, $I\,\ddot{\theta} = m g l\sin\theta + m l\cos\theta\,\ddot{x}$,
obtemos

\[
\boxed{\;\dot{E} = m l\cos\theta\,\dot{\theta}\,\ddot{x}\;}
\]

--- os termos gravitacionais se cancelam e a energia só varia por $\ddot{x}$.
Repetindo o argumento de Lyapunov com
$\dot{V} = \tilde{E}\,\dot{E} = \tilde{E}\,m l\cos\theta\,\dot{\theta}\,\ddot{x}$
e escolhendo $\ddot{x}\propto -\,\dot{\theta}\,\tilde{E}\cos\theta$, chegamos a
$\dot{V} = -\,k\,m l\,\tilde{E}^{2}\,\dot{\theta}^{2}\cos^{2}\theta \le 0$. A lei
de \emph{energy shaping} adaptada ao cart-pole fica, então, a de~\cite{mit} com
o fator geométrico $\cos\theta$:

\[
\boxed{\;u = -\,k\,\dot{\theta}\,\tilde{E}\,\cos\theta,
\qquad u \leftarrow \operatorname{sat}_{\pm u_{max}}(u).\;}
\]

Quando falta energia ($\tilde{E}<0$) a lei injeta energia
($\dot{E}>0$); ao atingir $E^{d}$ o termo se anula. Nos pontos de retorno
($\dot{\theta}=0$) e na horizontal ($\cos\theta=0$) tem-se $u=0$.

\subsection{Regulação da posição do carro}

A lei de \emph{energy shaping} pura tem uma limitação conhecida: governa a
energia da haste, mas não diz nada sobre a posição do carro. Como o
bombeamento empurra o carro sempre no sentido que adiciona energia, o carro
deriva sem limite (nas simulações, dezenas de metros), e o LQR recebe a haste
no topo com o carro deslocado --- fora da região de atração, e a captura
falha.

A solução clássica, seguindo Tedrake~\cite{mit}, é somar à lei de
\emph{energy shaping} um controlador PD sobre o carro, que o mantém próximo da
origem sem impedir o balanço:

\[
\boxed{\;u = -\,k_{swing}\,\dot{\theta}\,\tilde{E}\cos\theta
\;\underbrace{-\,k_p\,x\;-\;k_d\,\dot{x}}_{\text{PD no carro}},
\qquad u \leftarrow \operatorname{sat}_{\pm u_{max}}(u).\;}
\]

O termo $-k_p\,x$ age como uma mola que puxa o carro de volta a $x=0$ e
$-k_d\,\dot{x}$ como amortecimento. Optamos por um PD (e não PI/PID): o
\emph{swing-up} é um transitório curto sob uma ``perturbação'' oscilatória (o
próprio bombeamento troca de sinal a cada balanço), sem erro de regime
permanente a corrigir --- um termo integral só acumularia esse sinal
(\emph{windup}) e brigaria com o bombeamento. Os ganhos, obtidos por busca em
simulação, são

\[
k_{swing} = 12, \qquad k_p = 5, \qquad k_d = 2,
\]

com os quais a haste sobe e é capturada em torno de $1{,}2$~s (estado ideal,
sem ruído). É um ponto de partida: os ganhos podem ser reotimizados, por
exemplo minimizando o tempo de subida sob restrição da excursão do carro.

\subsection{Comutação para o LQR}

O \emph{swing-up} entrega a haste dentro da região de atração do LQR. Adotamos
uma condição de comutação por ângulo e velocidade pequenos,

\[
\lvert\theta\rvert < 15^{\circ}
\quad\text{e}\quad
\lvert\dot{\theta}\rvert < 2~\text{rad/s}
\;\Longrightarrow\;
\text{muda para } u = -K\hat{X},
\]

combinando a fase não-linear (energia) com a fase linear (LQR) num único
controlador de duas fases, exatamente o cenário que exercita os estimadores em
regimes distintos.

\subsection{Comportamento ideal de referência (sem ruído)}
\label{sec:ideal}

Antes de entrar no modelo de sensores, rodamos o cenário S0: medição perfeita
($\sigma_x=\sigma_\theta=0$, sem \emph{dropout} nem \emph{outliers}), para ver
o que o controlador entrega sozinho --- o alvo que os estimadores tentam
aproximar quando o ruído é reintroduzido. A Figura~\ref{fig:ideal} mostra
$\theta(t)$ e $x(t)$ para a mesma condição inicial de S2 (haste pendurada,
$\theta_0=180^{\circ}$): o \emph{swing-up} bombeia energia até a haste cruzar
a vizinhança de $15^{\circ}$/$2$~rad/s em torno da vertical, o controlador
comuta para o LQR, e a haste estabiliza em $\theta\to0$ com o carro
recentrado ($x\to0$) em poucos segundos.

\figplaceholder{results/S0\_ideal\_traj.png}{Comportamento ideal (S0, sem
ruído de medição): ângulo $\theta$ e posição do carro $x$ sob controle real
(\emph{swing-up} $\to$ LQR) a partir da haste pendurada. Com a medição exata,
o estado estimado coincide com a verdade (KF/EKF/UKF, Seção~\ref{sec:ideal-rmse}) e
esta é, portanto, a trajetória de malha fechada ``sem perdas'' — a referência
contra a qual os cenários com ruído (S1--S6) são comparados.}
{fig:ideal}

%-----------------------------------------------------------------------------------------------
%   MODELO DE SENSORES
%-----------------------------------------------------------------------------------------------
% >>> [SEÇÃO B] label adicionado só para permitir referência cruzada a partir da
% nova Seção "Estimadores de Estado Clássicos"; conteúdo desta seção não foi alterado.
\section{Modelo de sensores e cenários de teste}
\label{sec:sensores}

O modelo de sensores é a interface entre a ``verdade'' do simulador e o que os
estimadores observam: recebe o estado verdadeiro e devolve uma medição
imperfeita. Como os filtros nunca acessam $X$ diretamente, é esse modelo que
torna o problema de estimação realista.

\subsection{Grandezas medidas}

Medimos só as posições, reproduzindo um par de encoders na posição do carro e
no ângulo da haste; as velocidades $\dot{x}$ e $\dot{\theta}$ não são medidas
e precisam ser estimadas. A medição é, portanto,

\[
z_k =
\begin{bmatrix} x \\ \theta \end{bmatrix}_{t_k}
+\; v_k,
\qquad
v_k \sim \mathcal{N}(0,\,R),
\qquad
R = \operatorname{diag}\!\left(\sigma_x^{2},\ \sigma_\theta^{2}\right),
\]

com ruído branco gaussiano e independente por canal. Os desvios nominais são
$\sigma_x = \sigma_\theta = 0{,}005$ (Tabela~\ref{tab:params}). Medir só
posições é o que dá sentido à comparação de estimadores: obter $\dot{x}$ e
$\dot{\theta}$ é exatamente a tarefa que separa a diferenciação com
passa-baixas dos filtros baseados em modelo (KF/EKF/UKF).

\subsection{Degradações e escolhas de projeto}

Além do ruído nominal, o sensor reproduz três degradações, acionadas pelos
parâmetros dos cenários:

\begin{itemize}
  \item \textbf{Ruído elevado (S3).} Multiplicamos $\sigma_x$ e $\sigma_\theta$
  por $10\times$. Testa a filtragem quando a relação sinal-ruído cai,
  favorecendo estimadores que ponderam corretamente medição \emph{vs.} modelo.

  \item \textbf{\emph{Dropout} de medida (S4).} A cada passo, com probabilidade
  $p_{drop}$, a leitura é totalmente perdida. Sinalizamos a perda com uma
  medição vazia ($z_k=\varnothing$), em vez de um valor qualquer, para que o
  estimador execute um passo de predição pura (sem atualização): a média é
  propagada pelo modelo e a covariância $P$ cresce, refletindo o aumento de
  incerteza. Zerar ou congelar a medição daria ao filtro uma confiança falsa,
  corrompendo a análise de consistência (NEES/NIS).

  \item \textbf{\emph{Outliers}.} Cada canal, com probabilidade $p_{out}$,
  recebe um pico espúrio de amplitude proporcional a $\sigma$ (por padrão
  $30\sigma$) e sinal aleatório. Modela leituras absurdas e ocasionais, que um
  estimador robusto deve rejeitar. Fica desligado por padrão ($p_{out}=0$) e
  serve para um teste opcional de robustez.
\end{itemize}

Todas as degradações têm probabilidade nula por padrão, de forma que os
cenários sem \emph{override} usam só o ruído gaussiano nominal e ficam
comparáveis entre si.

%=================================================================================================
%   ESTIMADORES DE ESTADO CLÁSSICOS
%   Equações e valores numéricos conferidos contra o código (src/est_lowpass.m,
%   src/est_complementary.m, src/est_kf.m, src/params.m) e contra um cálculo
%   numérico independente (discretização ZOH).
%=================================================================================================
\section{Estimadores de Estado Clássicos: da Diferenciação ao Filtro de Kalman Linear}
\label{sec:estimadores-classicos}

Esta seção descreve os três estimadores "clássicos" do projeto --- os que não
dependem de linearização \emph{online} nem de amostragem estocástica do espaço
de estados, ao contrário do EKF e do UKF (Seção~\ref{sec:ekf-ukf}). Os três
compartilham a mesma interface de \emph{drop-in} do contrato do projeto,

\[
[\hat X_k,\ P_k] = \texttt{estimator\_update}(\hat X_{k-1},\, P_{k-1},\, z_k,\, u_k,\, p),
\]

e atuam sobre a mesma medição de posições $z_k=[x\ \ \theta]^{\mathsf T}_{t_k}$
(Seção~\ref{sec:sensores}), mas diferem na quantidade de modelo que carregam:
da ausência total de modelo dinâmico (baseline), passando por um modelo
implícito (complementar), até o modelo linear completo do pêndulo, propagado
com sua incerteza estatística (Kalman linear). Essa progressão prepara o
terreno para o EKF e o UKF (seção seguinte), que estendem o mesmo princípio de
propagação de covariância ao modelo não-linear do cart-pole.

\subsection{Baseline do curso: diferenciação numérica com filtro passa-baixas}
\label{sec:baseline-lowpass}

O primeiro estimador (\texttt{est\_lowpass.m}) reproduz exatamente a técnica
vista em CMC-12: um filtro passa-baixas digital de primeira ordem (média móvel
exponencial) aplicado à medição, seguido de diferenciação numérica para obter as
velocidades~\cite{oppenheim}. Para cada canal medido,

\[
\tilde x_k = \alpha\,\tilde x_{k-1} + (1-\alpha)\,z_{x,k},
\qquad
\tilde\theta_k = \alpha\,\tilde\theta_{k-1} + (1-\alpha)\,z_{\theta,k},
\qquad \alpha\in(0,1),
\]

e as velocidades são obtidas por diferença atrasada (\emph{backward
difference}) do sinal já filtrado:

\[
\hat{\dot x}_k = \frac{\tilde x_k - \tilde x_{k-1}}{\Delta t},
\qquad
\hat{\dot\theta}_k = \frac{\tilde\theta_k - \tilde\theta_{k-1}}{\Delta t}.
\]

Esse filtro recursivo de primeiro grau é o análogo discreto de um passa-baixas
$RC$ contínuo com constante de tempo $\tau = -\Delta t/\ln\alpha$, ou
frequência de corte
$f_c = \dfrac{1}{2\pi\tau} = \dfrac{-\ln\alpha}{2\pi\,\Delta t}$
\cite{oppenheim}. Com o valor padrão do código, $\alpha_{lp}=0{,}8$, e
$\Delta t = 0{,}01$~s (Tabela~\ref{tab:params}), $f_c \approx 3{,}55$~Hz:
qualquer conteúdo do sinal acima dessa frequência --- inclusive parte da
dinâmica real do \emph{swing-up}, rápida e fortemente não-linear --- é
atenuado e atrasado junto com o ruído do encoder.

Esse é o ponto fraco do baseline, e a razão de ele servir de referência
inferior na comparação: não há nenhum modelo da planta embutido no filtro.
Quanto maior $\alpha$, menor o ruído mas maior o atraso de fase --- o mesmo
compromisso ruído \emph{vs.} atraso de qualquer passa-baixas clássico
\cite{oppenheim} --- e não dá para reduzir os dois ao mesmo tempo, porque o
filtro não distingue ruído de medição de dinâmica real. A covariância $P$
também não tem significado probabilístico aqui: o código só a repassa adiante
(\texttt{P = P\_prev}) para respeitar a interface comum, já que não há modelo
de incerteza a propagar.

\subsection{Filtro complementar}
\label{sec:complementar}

O segundo estimador (\texttt{est\_complementary.m}) já incorpora uma noção
mínima de modelo: a cada passo, prevê a posição por integração de Euler à
frente supondo velocidade constante,

\[
x_{pred} = \hat x_{k-1} + \Delta t\,\hat{\dot x}_{k-1},
\qquad
\theta_{pred} = \hat\theta_{k-1} + \Delta t\,\hat{\dot\theta}_{k-1},
\]

e funde essa previsão com a medição bruta usando um ganho fixo
$\alpha_{comp}\in(0,1)$,

\[
x_{fused} = (1-\alpha_{comp})\,x_{pred} + \alpha_{comp}\,z_{x,k},
\qquad
\theta_{fused} = (1-\alpha_{comp})\,\theta_{pred} + \alpha_{comp}\,z_{\theta,k},
\]

com as velocidades novamente obtidas por diferença atrasada do sinal fundido.
O nome "complementar" vem da interpretação clássica no domínio da frequência:
dados dois sensores (ou, aqui, um
sensor e um modelo) que medem a mesma grandeza com espectros de ruído
complementares, pondera-se cada fonte por um par de filtros $H(s)$ e $1-H(s)$
que somam a unidade —

\[
\hat X(s) = H(s)\,Z(s) + \big(1-H(s)\big)\,X_{pred}(s),
\]

de modo que a medição (confiável em baixa frequência, mas ruidosa em alta) e o
modelo/previsão (suave, mas sujeito a deriva em regime permanente) se
complementam sem precisar estimar covariâncias explicitamente. A versão
discreta a ganho fixo implementada aqui é a forma mais simples dessa ideia.
Com o ganho padrão do código, $\alpha_{comp}=0{,}02$, o filtro confia
fortemente na previsão (98\,\%) e corrige lentamente pela medição (2\,\%) a
cada passo --- adequado quando o modelo de velocidade constante é razoável
entre amostras, mas ainda assim acumula deriva se ficar muitos passos sem
medição, o problema clássico de fusão de sensores~\cite{simon2006}. É
justamente esse ganho fixo e invariante no tempo que o Filtro de Kalman
generaliza na próxima subseção: o KF pode ser visto como um filtro
complementar cujo ganho $K_k$ é recalculado a cada passo para ser o ótimo
dado o par modelo/medição~\cite{simon2006}.

\subsection{Do modelo linearizado ao Filtro de Kalman discreto}
\label{sec:kf-linear}

\subsubsection{Discretização por segurador de ordem zero (ZOH)}
\label{sec:zoh}

O terceiro estimador (\texttt{est\_kf.m}) usa o modelo linear $\dot X = AX+Bu$
já obtido na Seção~\ref{sec:linearizacao} em torno da vertical, complementado
pelo modelo de medição (linear e exato, já que só observamos posições
diretamente),

\[
z = C X, \qquad
C = \begin{bmatrix} 1 & 0 & 0 & 0 \\ 0 & 0 & 1 & 0 \end{bmatrix}.
\]

Como o KF discreto opera sobre o par $(A_d, B_d)$ amostrado a $\Delta t$, e não
sobre $(A,B)$ contínuos, precisamos discretizar. A discretização exata de um
sistema linear sob entrada constante por trecho (\emph{zero-order hold}, ZOH) é

\[
A_d = e^{A\Delta t}, \qquad
B_d = \int_0^{\Delta t} e^{A\tau}\,d\tau\;B,
\]

resultado padrão de controle digital~\cite{simon2006}.
Quando $A$ é invertível, $B_d = A^{-1}(A_d - I)B$; não é o caso aqui: a
dinâmica do cart-pole é invariante à translação (a posição $x$ não aparece em
nenhuma equação de movimento), a primeira coluna de $A$
(Seção~\ref{sec:linearizacao}) é inteiramente nula e $A$ é singular. O caminho
numericamente robusto --- e o que o MATLAB usa internamente em \texttt{c2d} ---
é formar a exponencial da matriz aumentada~\cite{simon2006}

\[
\exp\!\left(
\begin{bmatrix} A & B \\ 0 & 0 \end{bmatrix}\Delta t
\right)
=
\begin{bmatrix} A_d & B_d \\ 0 & I \end{bmatrix},
\]

que evita a inversão e vale mesmo com $A$ singular. Como a saída $z=CX$ é uma
relação puramente algébrica (sem dinâmica), a discretização não altera $C$:
$C_d = C$. É esse procedimento que a função auxiliar
\texttt{linearise\_upright.m} executa via
\texttt{c2d(ss(A\_c,B\_c,C\_c,0), dt, `zoh')}, reaproveitando as matrizes $A$ e
$B$ já derivadas e verificadas na Seção~\ref{sec:linearizacao}; o resultado
($p.Ad$, $p.Bd$, $p.Cd$) é calculado uma única vez dentro de \texttt{params.m}
e fica disponível a \texttt{est\_kf.m} pelo struct de parâmetros. Conferimos a
consistência entre essa linearização e o Jacobiano analítico de
\texttt{jacobian\_f.m} (Seção~\ref{sec:ekf-ukf}, derivado independentemente
para o EKF) em \texttt{tests/test\_classical\_filters.m}: as duas
discretizações batem a menos de $10^{-8}$ (Seção~\ref{sec:sintese-classicos}).

Com $\Delta t = 0{,}01$~s e os valores numéricos de $A$ e $B$ já apresentados
na Seção~\ref{sec:linearizacao}, a discretização ZOH dá (calculado a partir da
matriz aumentada acima)

\[
A_d \approx
\begin{bmatrix}
1 & 0{,}009995 & 1{,}03\times10^{-4} & 3{,}4\times10^{-7} \\
0 & 0{,}999070 & 0{,}020529          & 1{,}03\times10^{-4} \\
0 & -6{,}98\times10^{-6} & 1{,}000890 & 0{,}010003 \\
0 & -0{,}001395 & 0{,}177987          & 1{,}000890
\end{bmatrix},
\qquad
B_d \approx
\begin{bmatrix}
4{,}65\times10^{-5} \\ 0{,}009298 \\ 6{,}98\times10^{-5} \\ 0{,}013951
\end{bmatrix}.
\]

Como esperado de uma discretização fiel a um passo pequeno frente à dinâmica
($\Delta t=0{,}01$~s $\ll$ constantes de tempo de $A$, polos em $\pm4{,}22$ e
$-0{,}077$~rad/s --- Seção~\ref{sec:linearizacao}), $A_d$ fica próxima da
identidade e os autovalores de $A_d$ (isto é, $e^{\lambda_i(A)\Delta t}$) ficam
próximos de $1$: um deles, $e^{+4{,}2266\times0{,}01}\approx1{,}0430$, maior
que $1$ em módulo, confirma que o modo instável da haste permanece instável
também no domínio discreto, como esperado.

\subsubsection{Algoritmo do Filtro de Kalman}
\label{sec:kf-algoritmo}

Com $(A_d,B_d,C_d)$ em mãos, o KF linear implementado segue a forma padrão de
predição--correção de Kalman~\cite{welch1995, simon2006}:

\medskip
\noindent\textbf{Predição} (propaga estado e incerteza pelo modelo, sem usar a medição):
\[
\hat X_{k}^{-} = A_d\,\hat X_{k-1} + B_d\,u_{k-1},
\qquad
P_{k}^{-} = A_d\,P_{k-1}\,A_d^{\mathsf T} + Q,
\]

\noindent\textbf{Atualização} (corrige com a nova medição $z_k$):
\[
S_k = C_d\,P_k^{-}\,C_d^{\mathsf T} + R,
\qquad
K_k = P_k^{-}\,C_d^{\mathsf T}\,S_k^{-1},
\qquad
\hat X_k = \hat X_k^{-} + K_k\big(z_k - C_d\hat X_k^{-}\big),
\qquad
P_k = (I - K_k C_d)\,P_k^{-},
\]

em que $Q$ e $R$ são as covariâncias (supostas conhecidas, diagonais e
constantes) do ruído de processo e de medição (Tabela~\ref{tab:params-qr}
adiante). O ganho $K_k$ não é arbitrário: é o único ganho que minimiza a
incerteza remanescente $\operatorname{tr}(P_k)$ a cada passo. Partindo da
forma de Joseph
$P_k=(I-K_kC_d)P_k^{-}(I-K_kC_d)^{\mathsf T}+K_kRK_k^{\mathsf T}$ (numericamente
mais estável, adotada também no EKF do projeto), derivamos
$\operatorname{tr}(P_k)$ em relação a $K_k$ e igualamos a zero, obtendo

\[
\frac{\partial}{\partial K_k}\operatorname{tr}(P_k) = 0
\;\Longrightarrow\;
K_k = P_k^{-}C_d^{\mathsf T}\big(C_dP_k^{-}C_d^{\mathsf T}+R\big)^{-1} = P_k^{-}C_d^{\mathsf T}S_k^{-1},
\]

que é exatamente o ganho acima~\cite{simon2006}. Sob as hipóteses usuais ---
modelo linear, ruídos brancos, gaussianos, de média nula e descorrelacionados
entre si e do estado inicial --- esse ganho torna o KF o estimador de
variância mínima entre todos os estimadores lineares não enviesados
(propriedade BLUE), e coincide com o estimador de máxima verossimilhança/MMSE
quando os ruídos são gaussianos~\cite{simon2006}. É essa otimalidade
comprovada --- e não uma escolha heurística de ganho, como no filtro
complementar (Seção~\ref{sec:complementar}) --- que justifica o KF como o
patamar de comparação "ótimo" dentro do regime linear para o restante do
trabalho.

\paragraph{Escolha de $Q$ e $R$.} O código (\texttt{params.m}) adota como
ponto de partida

\begin{table}[H]
\centering
\begin{tabular}{l l}
\hline
\textbf{Matriz} & \textbf{Valor} \\
\hline
$R$ (ruído de medição)  & $\operatorname{diag}(\sigma_x^2,\ \sigma_\theta^2) = \operatorname{diag}(2{,}5\times10^{-5},\ 2{,}5\times10^{-5})$ \\
$Q$ (ruído de processo) & $\operatorname{diag}(10^{-4},\ 10^{-3},\ 10^{-4},\ 10^{-3})$ \\
\hline
\end{tabular}
\caption{Covariâncias de ruído usadas pelo KF linear (compartilhadas com EKF/UKF
via \texttt{p.Q}/\texttt{p.R}, salvo \emph{override} \texttt{p.Q\_kf}/\texttt{p.R\_kf}).}
\label{tab:params-qr}
\end{table}

$R$ decorre direto do ruído físico dos encoders (Tabela~\ref{tab:params}) e
não é, a rigor, um parâmetro de sintonia. Já $Q$ modela a incerteza sobre o
quanto o modelo linearizado deixa de representar a dinâmica real (erro de
linearização, atrito não modelado, discretização) e na prática é um parâmetro
de ajuste: $Q$ pequeno demais faz o filtro confiar demais no modelo linear e
divergir fora da vizinhança da vertical (o regime do \emph{swing-up}); $Q$
grande demais aproxima o KF de uma fusão ponderada simples, perdendo a
vantagem de usar o modelo. Ajustamos $Q$ diretamente no tempo discreto, e não
pela discretização rigorosa de uma densidade espectral de ruído de processo
contínuo --- simplificação comum em projetos aplicados~\cite{simon2006}, mas
que vale registrar: $Q$ aqui é "efetivo", não "físico".

A consistência estatística desse ajuste de $Q$ é o que o teste NEES
(\emph{Normalised Estimation Error Squared}), em \texttt{compute\_metrics.m},
mede: sob as hipóteses acima, $e_k^{\mathsf T}P_k^{-1}e_k$ segue uma
distribuição $\chi^2$ com 4 graus de liberdade (dimensão do estado); NEES
sistematicamente acima do intervalo esperado indica $P$ otimista (subestima o
erro real, tipicamente $Q$ ou $R$ pequenos demais), e abaixo indica $P$
pessimista (superestima o erro)~\cite{simon2006}. O mesmo raciocínio, aplicado
à inovação $\nu_k=z_k-C_d\hat X_k^{-}$ em vez do erro de estado, dá o teste
NIS ($\chi^2$ com 2 graus de liberdade), também em \texttt{compute\_metrics.m}.

\paragraph{Sensibilidade a $Q/R$ mal ajustados (cenário S5).} O cenário S5
testa exatamente esse ponto: o sensor continua gerando ruído nominal, mas os
estimadores recebem covariâncias erradas, $Q_{\text{visto}}=10\,Q$ e
$R_{\text{visto}}=R/10$ (Seção~\ref{sec:ekf-ukf} traz a tabela completa com os
cinco estimadores). O efeito é o previsto pela teoria: o KF confia mais na
medição e menos na predição do que deveria, subestimando o erro real que sua
própria covariância indica. O NEES médio do KF sobe de $1{,}66$ em S1 (dentro
do intervalo $\chi^2_4$ a 95\%, $[0{,}48;\,11{,}14]$) para $19{,}19$ em S5 ---
acima do limite superior, confirmando $P$ otimista. O RMSE de $\theta$ quase
não piora ($0{,}24^{\circ}\to0{,}29^{\circ}$): o ganho ainda funciona
razoavelmente bem na prática, mas o filtro passa a relatar uma incerteza
menor que a real --- o tipo de falha que só o teste de consistência expõe, e
que um RMSE sozinho esconderia.

\subsection{Síntese comparativa}
\label{sec:sintese-classicos}

A Tabela~\ref{tab:classicos} resume as diferenças estruturais entre os três
estimadores desta seção.

% >>> [SEÇÃO B] Tabela 3 reformatada: estava com 4 colunas "l" (largura livre)
% e estourava a margem da página porque o cabeçalho da última coluna era uma
% frase longa. Troquei para colunas de largura fixa (p{...}, quebram linha) e
% fonte \small, e encurtei o cabeçalho da última coluna.
\begin{table}[H]
\centering
\small
\begin{tabular}{@{}p{3.1cm} p{3.4cm} p{2.6cm} p{3.3cm}@{}}
\hline
\textbf{Estimador} & \textbf{Modelo dinâmico} & \textbf{Ganho} & \textbf{$P$ tem significado estatístico?} \\
\hline
Baseline (passa-baixas)  & Nenhum                        & Fixo ($\alpha_{lp}$)             & Não \\[2pt]
Complementar             & Implícito (veloc. constante)  & Fixo ($\alpha_{comp}$)           & Não \\[2pt]
KF linear                & $A_d,B_d$ (linearizado)       & Ótimo, variável ($K_k$)          & Sim (sob hipóteses lineares/gaussianas) \\
\hline
\end{tabular}
\caption{Comparação estrutural dos três estimadores clássicos implementados.}
\label{tab:classicos}
\end{table}

\paragraph{Metodologia.} Antes de avaliar os cinco estimadores na malha
fechada completa e nos seis cenários S1--S6 (Seção~\ref{sec:ekf-ukf}), isolamos
os três estimadores clássicos num teste controlado
(\texttt{tests/test\_classical\_filters.m}), que (i) confere numericamente a
linearização de \texttt{linearise\_upright.m} contra o Jacobiano analítico
independente \texttt{jacobian\_f.m} --- as duas batem a menos de $10^{-8}$;
(ii) simula uma trajetória próxima ao cenário S1 (haste a $5^{\circ}$ da
vertical) sob realimentação do estado verdadeiro ($u=-K_{demo}X$), com o ganho
$K$ já publicado na Seção~\ref{sec:matlab} --- escolha deliberada para isolar
o desempenho de cada estimador dos efeitos do controlador, retomada na
discussão abaixo; e (iii) roda os três estimadores em paralelo sobre essa
trajetória, com ruído nominal e com o multiplicador $10\times$ do cenário S3.

\paragraph{Resultados.} A Tabela~\ref{tab:rmse-classicos} e as
Figuras~\ref{fig:s1-classicos-theta}--\ref{fig:classicos-rmse} resumem o que
observamos.

\begin{table}[H]
\centering
\small
\begin{tabular}{@{}l r r r r@{}}
\hline
\textbf{Estimador} & \textbf{RMSE $\theta$, S1 (°)} & \textbf{RMSE $\theta$, S3 (°)} & \textbf{RMSE $x$, S1 (m)} & \textbf{RMSE $x$, S3 (m)} \\
\hline
Baseline (passa-baixas) & 0,250 & 1,018 & 0,0031 & 0,0175 \\
Complementar             & 1,036 & 2,500 & 0,0093 & 0,0388 \\
KF linear                & 0,241 & 1,157 & 0,0043 & 0,0188 \\
\hline
\end{tabular}
\caption{RMSE de $\theta$ e de $x$ para os três estimadores clássicos, no
cenário de demonstração descrito acima, com ruído nominal (S1) e com
ruído $10\times$ (S3). Valores gerados por
\texttt{tests/test\_classical\_filters.m}.}
\label{tab:rmse-classicos}
\end{table}

Dois pontos merecem destaque --- um deles contraria a expectativa ingênua de
que "mais modelo = sempre melhor":

\begin{itemize}
  \item \textbf{Baseline e KF praticamente empatados, ambos à frente do
  complementar.} No regime quase-linear e bem amostrado ($100$~Hz) deste
  cenário, o passa-baixas bem ajustado ($\alpha_{lp}=0{,}8$) e o KF linear
  entregam RMSE de $\theta$ quase idênticos ($0{,}25^{\circ}$ \emph{vs.}
  $0{,}24^{\circ}$ com ruído nominal). Isso não contradiz a otimalidade do KF
  (Seção~\ref{sec:kf-algoritmo}): o KF é ótimo dado o par $(Q,R)$
  configurado, e nesse cenário a trajetória verdadeira vem de uma EDO
  determinística, sem ruído de processo de fato injetado --- ou seja, o $Q$
  de \texttt{params.m} (Tabela~\ref{tab:params-qr}), calibrado para
  representar erro de linearização e dinâmica não modelada, fica relativamente
  "generoso" frente ao ruído de processo real (próximo de zero) deste cenário.
  Isso dá ao ganho do KF mais folga do que o necessário, deixando passar um
  pouco mais de ruído de medição do que um passa-baixas fortemente suavizante
  deixaria --- o mesmo fenômeno de sensibilidade a $Q/R$ mistunado que o
  cenário S5 expõe (Seção~\ref{sec:kf-algoritmo}). A vantagem estrutural do KF
  (explorar $A_d,B_d$) deve ficar mais evidente com ruído de processo real ou
  com $Q$ recalibrado --- ambos fora do escopo desta seção.

  \item \textbf{Filtro complementar é o pior dos três, por larga margem.} A
  Figura~\ref{fig:s1-classicos-theta} mostra o motivo: com o ganho padrão
  $\alpha_{comp}=0{,}02$ (Seção~\ref{sec:complementar}), o filtro confia
  $98\,\%$ na predição de velocidade constante a cada passo e corrige devagar
  pela medição --- adequado para deriva lenta, mas lento demais para
  acompanhar a oscilação inicial rápida de $5^{\circ}$ deste cenário: o filtro
  passa de $6^{\circ}$ e só converge por volta de $t\approx2$~s, enquanto
  baseline e KF acompanham a verdade quase de imediato. Sob ruído $10\times$
  (S3), seu RMSE mais que dobra ($1{,}04^{\circ}\to2{,}50^{\circ}$), a maior
  degradação relativa das três --- esperado, já que $\alpha_{comp}$ é fixo e
  não se adapta ao novo nível de ruído, ao contrário do ganho $K_k$ do KF,
  recalculado a cada passo.
\end{itemize}

Não testamos os três estimadores desta seção no \emph{swing-up} (cenário S2,
ângulos grandes): a comparação acima usa de propósito a realimentação do
estado verdadeiro ($u=-K_{demo}X$), e não da estimativa, para isolar o
desempenho dos filtros dos efeitos do controlador antes de avaliar a malha
fechada por completo. É justamente no swing-up
que a linearização por trás do KF deixa de valer; a comparação dos cinco
estimadores nesse regime, com a malha fechada real ($u=-K\hat X$ via
\texttt{swingup.m}), está na Seção~\ref{sec:ekf-ukf}.

\figplaceholder{results/S1\_classicos\_theta.png}{Ângulo $\theta$ estimado
pelos três estimadores clássicos no cenário de demonstração (análogo a S1,
$5^{\circ}$ de perturbação inicial, ruído nominal) \emph{vs.} estado verdadeiro
simulado. O filtro complementar (ganho fixo $\alpha_{comp}=0{,}02$)
\emph{overshoot}a e leva ${\sim}2$~s para convergir; baseline e KF acompanham a
verdade quase imediatamente.}{fig:s1-classicos-theta}

\figplaceholder{results/S1\_S3\_classicos\_rmse.png}{RMSE de $\theta$ dos três
estimadores clássicos no cenário de demonstração, com ruído nominal e com
ruído $10\times$ (análogo a S1 \emph{vs.} S3). O complementar tem a maior
degradação relativa; baseline e KF permanecem próximos em ambos os
níveis.}{fig:classicos-rmse}

\paragraph{Da comparação isolada à malha fechada completa.} A planta
não-linear (\texttt{plant\_cartpole.m}), o sensor com
\emph{dropout}/\emph{outliers} (\texttt{sensor\_model.m}), o LQR
(\texttt{controller\_lqr.m}) e o \emph{swing-up} (\texttt{swingup.m}) são
integrados por \texttt{run\_scenario.m}/\texttt{main.m} (\texttt{src/main.m}
roda os cinco estimadores nos seis cenários S1--S6 e salva
\texttt{results/*.mat}). A Seção~\ref{sec:ekf-ukf} traz os resultados dessa
malha real --- com uma diferença importante frente à comparação isolada
acima: ali o ganho $K_{demo}$ realimentava o estado verdadeiro; na malha real
$u=-K\hat X$ realimenta a estimativa, então um estimador ruim degrada o
próprio controle (e não só o RMSE do estimador isolado) --- o princípio da
separação em ação, não só em teoria.

%=================================================================================================
%   ESTIMADORES NÃO-LINEARES (EKF/UKF)
%   Resultados gerados pelo harness completo (src/main.m) sobre os seis
%   cenários S1..S6, via tests/build_report_results.m, a partir de
%   results/*.mat.
%=================================================================================================
\section{Estimadores de Estado Não-Lineares: EKF e UKF}
\label{sec:ekf-ukf}

O KF linear (Seção~\ref{sec:kf-linear}) só é ótimo perto do ponto de
linearização $\theta=0$; fora dessa vizinhança --- o regime do
\emph{swing-up} (cenário S2) --- nada garante nem sua estabilidade. Esta
seção trata os dois estimadores que lidam com a não-linearidade diretamente:
o Filtro de Kalman Estendido (EKF), que relineariza a cada passo em torno da
estimativa corrente, e o Filtro de Kalman Unscented (UKF), que propaga a
incerteza por um conjunto determinístico de pontos (\emph{sigma points})
através do modelo não-linear exato, sem linearização analítica. Ambos seguem
a mesma interface \emph{drop-in} das seções anteriores,
$[\hat X_k,\,P_k] = \texttt{estimator\_update}(\hat X_{k-1},\,P_{k-1},\,z_k,\,u_k,\,p)$,
e retornam opcionalmente $\texttt{dbg} = (\nu_k, S_k)$ para o teste de
consistência NIS em \texttt{compute\_metrics.m}.

\subsection{Filtro de Kalman Estendido (EKF)}
\label{sec:ekf}

\subsubsection{Jacobiano analítico}

O EKF precisa, a cada passo, de uma aproximação linear local da dinâmica em
torno da estimativa corrente $\hat X_{k-1}$. Em vez de linearizar uma única
vez (Seção~\ref{sec:linearizacao}), \texttt{jacobian\_f.m} calcula
$F=\partial f/\partial X$ analiticamente, em qualquer estado, reaproveitando a
mesma solução do sistema $2\times2$ de \texttt{plant\_cartpole.m}
($\Delta=aJ-h^2$, Seção~\ref{sec:linearizacao}) e aplicando a regra do
quociente às expressões de $\ddot x$ e $\ddot\theta$. Como a dinâmica não
depende de $x$, a coluna correspondente de $F$ é sempre nula; o restante é
avaliado em $(\hat x_2,\hat x_3,\hat x_4,u)$ a cada chamada. Verificamos a
corretude contra diferenças finitas centrais e contra a linearização de
\texttt{linearise\_upright.m} em $\theta=0$ (Seção~\ref{sec:ekf-verificacao}).

\subsubsection{Predição: integração não-linear + covariância linearizada localmente}

Diferentemente do KF, cuja predição é uma multiplicação matricial fixa
($A_d\hat X_{k-1}+B_d u_{k-1}$), o EKF integra a dinâmica não-linear completa
(RK4, mesmo integrador da planta verdadeira) para propagar a média:

\[
\hat X_k^{-} = \mathrm{RK4}\big(f,\ \hat X_{k-1},\ u_{k-1},\ \Delta t\big),
\]

e usa o Jacobiano \emph{congelado} em $\hat X_{k-1}$, discretizado pela mesma
exponencial de matriz da Seção~\ref{sec:zoh} (sem termo de entrada, pois só a
covariância é propagada dessa forma):

\[
F = \texttt{jacobian\_f}(\hat X_{k-1}, u_{k-1}, p),
\qquad
F_d = e^{F\Delta t},
\qquad
P_k^{-} = F_d\,P_{k-1}\,F_d^{\mathsf T} + Q.
\]

A média segue a física exata; só a incerteza usa uma aproximação linear
local, recalculada a cada passo no ponto atual (ao contrário do $(A_d,B_d)$
fixo do KF) --- é essa relinearização contínua que permite ao EKF acompanhar
o \emph{swing-up} sem perder a estimativa
(Seção~\ref{sec:resultados-naolineares}).

\subsubsection{Atualização}

O modelo de medição continua linear ($z=CX$, mesma matriz $C$ da
Seção~\ref{sec:kf-linear}), de modo que a atualização tem a mesma forma do
KF, na forma de Joseph (Seção~\ref{sec:kf-algoritmo}):

\[
S_k = C P_k^{-} C^{\mathsf T} + R,
\qquad
K_k = P_k^{-} C^{\mathsf T} S_k^{-1},
\qquad
\hat X_k = \hat X_k^{-} + K_k\,\nu_k,
\qquad
\nu_k = z_k - C\hat X_k^{-},
\]
\[
P_k = (I-K_kC)\,P_k^{-}(I-K_kC)^{\mathsf T} + K_k R K_k^{\mathsf T}.
\]

Em \emph{dropout} (S4, $z_k=[\mathrm{NaN},\mathrm{NaN}]^{\mathsf T}$), o EKF
pula a atualização e retorna $(\hat X_k^{-}, P_k^{-})$, a mesma convenção dos
demais estimadores (Seção~\ref{sec:sensores}).

\subsection{Filtro de Kalman Unscented (UKF)}
\label{sec:ukf}

\subsubsection{Motivação: aproximar a distribuição, não a função}

O EKF lineariza a função de transição a cada passo --- uma aproximação de
primeira ordem que descarta termos de ordem $\geq2$ e introduz viés quando a
curvatura de $f$ é grande. O UKF~\cite{wan2000} aproxima, em vez disso, a
distribuição de probabilidade: escolhemos deterministicamente um pequeno
conjunto de pontos (\emph{sigma points}) que reproduzem exatamente a média e a
covariância de $\hat X_{k-1}$, propagamos cada um pela dinâmica não-linear
exata (RK4, sem aproximação) e reconstruímos a média e a covariância a
posteriori a partir dos pontos transformados. O resultado é exato até 2ª
ordem de Taylor da média (3ª se a entrada for gaussiana), uma ordem a mais que
o EKF, sem exigir Jacobiano.

\subsubsection{Sigma points e pesos (transformada \emph{unscented} escalonada)}

A implementação (\texttt{est\_ukf.m}) segue a forma escalonada de Van der
Merwe~\cite{wan2000}. Para $n=4$, três parâmetros controlam a dispersão:
$\alpha$ (tipicamente $10^{-3}$), $\beta$ ($=2$, ótimo para gaussianas) e
$\kappa$ (geralmente $0$), com

\[
\lambda = \alpha^2(n+\kappa) - n.
\]

Os $2n+1=9$ \emph{sigma points} são a média mais/menos as colunas da raiz de
Cholesky da covariância escalada,

\[
\mathcal{X}_0 = \hat X_{k-1},
\qquad
\mathcal{X}_i = \hat X_{k-1} \pm \Big(\sqrt{(n+\lambda)P_{k-1}}\Big)_i,
\quad i=1,\dots,n,
\]

com pesos

\[
W_0^{(m)} = \frac{\lambda}{n+\lambda}, \qquad
W_0^{(c)} = \frac{\lambda}{n+\lambda} + (1-\alpha^2+\beta), \qquad
W_i^{(m)} = W_i^{(c)} = \frac{1}{2(n+\lambda)}\ \ (i\geq1).
\]

Se a fatoração de Cholesky de $P_{k-1}$ falha (S6), o código regulariza
somando $10^{-6}I$ antes de tentar de novo.

\subsubsection{Predição}

Cada um dos $2n+1$ \emph{sigma points} é propagado pela dinâmica não-linear
completa (RK4, sem aproximação), e a média e a covariância previstas são
reconstruídas por combinação ponderada:

\[
\mathcal{X}_i^{-} = \mathrm{RK4}(f,\ \mathcal{X}_i,\ u_{k-1},\ \Delta t),
\qquad
\hat X_k^{-} = \sum_{i=0}^{2n} W_i^{(m)}\,\mathcal{X}_i^{-},
\]
\[
P_k^{-} = Q + \sum_{i=0}^{2n} W_i^{(c)}\big(\mathcal{X}_i^{-}-\hat X_k^{-}\big)\big(\mathcal{X}_i^{-}-\hat X_k^{-}\big)^{\mathsf T}.
\]

\subsubsection{Atualização}

Em geral o UKF propagaria os \emph{sigma points} também pelo modelo de medição
$h(\cdot)$. Aqui, porém, $h(X)=CX$ é linear (mesma matriz $C$ do KF/EKF), e
propagar os pontos por $h$ daria exatamente $C P_k^{-} C^{\mathsf T}+R$ --- o
mesmo resultado da forma fechada do KF. A implementação explora essa
igualdade direto, sem reamostrar:

\[
S_k = C P_k^{-} C^{\mathsf T} + R, \qquad
K_k = P_k^{-} C^{\mathsf T} S_k^{-1}, \qquad
\hat X_k = \hat X_k^{-} + K_k\,\nu_k, \qquad
\nu_k = z_k - C\hat X_k^{-},
\]
\[
P_k = (I-K_kC)\,P_k^{-}.
\]

Diferentemente do KF e do EKF, essa atualização não usa a forma de Joseph ---
como a álgebra é exata ($h$ já é linear), o ganho de estabilidade numérica é
menor aqui. O caso de \emph{dropout} segue a mesma convenção dos demais
estimadores: sem atualização, retornamos $(\hat X_k^{-},P_k^{-})$.

\subsection{Verificação numérica}
\label{sec:ekf-verificacao}

\texttt{tests/test\_nonlinear\_filters.m} valida três propriedades:

\begin{enumerate}
  \item \textbf{Jacobiano analítico \emph{vs.} diferença finita central.}
  Para seis valores de $\theta$ (incluindo $0$, $\pm\pi/2$ e $\pi$), o erro
  máximo entre \texttt{jacobian\_f.m} e diferença finita central de
  \texttt{plant\_cartpole.m} foi $2{,}51\times10^{-9}$, abaixo da tolerância
  de $10^{-6}$ em todo o domínio de $\theta$.
  \item \textbf{\emph{Smoke tests}.} Em duas trajetórias sintéticas, $\hat X$
  e $P$ ficam finitos e $P$ simétrica positiva-semidefinida (via
  \texttt{chol}) em todos os passos, mesmo atravessando \emph{dropout}, sem
  produzir \texttt{NaN}/\texttt{Inf}.
  \item \textbf{Interface.} Chamadas com e sem o argumento de depuração
  \texttt{dbg} funcionam; \texttt{dbg} vem vazio quando a medição é perdida.
\end{enumerate}

Os sete testes passam (\texttt{ALL TESTS PASSED}).

\subsection{Resultados: divergência do KF fora do regime linear (\emph{swing-up}, S2)}
\label{sec:resultados-naolineares}

Com a malha fechada completa (Seção~\ref{sec:sintese-classicos}), rodamos os
cinco estimadores no cenário mais exigente: o \emph{swing-up} (S2), em que a
haste passa por $\theta\in(-\pi,\pi]$ antes de ser capturada pelo LQR ---
exatamente o regime para o qual a linearização do KF em $\theta=0$ não foi
projetada.

\figplaceholder{results/S2\_swingup\_theta.png}{Ângulo $\theta$ verdadeiro
\emph{vs.} estimado pelos cinco estimadores durante o \emph{swing-up} completo
(S2, malha fechada real com \texttt{swingup.m}). Baseline e KF acompanham a
verdade razoavelmente bem enquanto o pêndulo ainda executa poucos balanços,
mas o KF passa a oscilar sem relação com a verdade assim que o ângulo cruza
regiões distantes de $\theta=0$; EKF e UKF (roxo/verde) colam na verdade
(preto) do início ao fim, inclusive durante os balanços de $180^{\circ}$.}
{fig:s2-swingup-theta}

\figplaceholder{results/S2\_nees.png}{NEES ao longo do \emph{swing-up} (S2)
para KF, EKF e UKF (escala log). A faixa tracejada é o intervalo
$\chi^2_4$ a 95\%, $[0{,}48;\,11{,}14]$. O KF passa a maior parte do tempo
\textbf{quatro a cinco ordens de grandeza} acima do limite superior — o
próprio filtro "acredita" ter convergido (sua covariância $P$ encolhe) ao
mesmo tempo em que o erro real explode, o padrão clássico de
\textbf{divergência otimista} de um filtro linear fora de sua região de
validade. EKF e UKF permanecem dentro (ou muito perto) da faixa esperada o
tempo todo.}{fig:s2-nees}

A Tabela~\ref{tab:s2-swingup} quantifica o que as Figuras~\ref{fig:s2-swingup-theta}
e~\ref{fig:s2-nees} mostram.

\begin{table}[H]
\centering
\small
\begin{tabular}{@{}l r r r r@{}}
\hline
\textbf{Estimador} & \textbf{RMSE $\theta$ (°)} & \textbf{RMSE $\dot\theta$ (rad/s)} & \textbf{NEES médio} & \textbf{NEES $\leq\chi^2_{4,97.5\%}$?} \\
\hline
Baseline      &  9,75 &  2,73 &        7,81 & aprox.\ (limítrofe) \\
Complementar  & 51,62 & 12,89 &      223,17 & não \\
KF linear     & 10,51 & 90,44 & 251699,13   & não (diverge) \\
EKF           &  0,24 &  0,04 &        1,73 & sim \\
UKF           &  0,23 &  0,03 &        1,58 & sim \\
\hline
\end{tabular}
\caption{Desempenho no \emph{swing-up} (S2), malha fechada real
($u=-K\hat X$ via \texttt{swingup.m}), Monte Carlo de uma semente (\texttt{rng}
fixado pelo harness). O RMSE de $\dot\theta$ do KF ($90{,}4$~rad/s) é o
sintoma mais direto da divergência: o ganho calculado a partir de um modelo
linearizado errado deixa de fazer sentido físico.}
\label{tab:s2-swingup}
\end{table}

O KF linear não é só pior: ele diverge de forma otimista --- sua própria
covariância $P$ relata confiança crescente exatamente quando o erro real
explode (Figura~\ref{fig:s2-nees}), o que o tornaria perigoso numa malha real.
O baseline, sem modelo algum a "proteger", simplesmente segue a medição
filtrada --- ruim (RMSE $9{,}75^{\circ}$) mas nunca catastrófico. Confirma-se
o argumento da Introdução: o KF linear compensa só no regime quase-linear;
fora dele, um modelo errado pode ser pior que nenhum modelo, e é aí que EKF e
UKF, com o modelo não-linear correto, se destacam.

\subsection{Síntese comparativa final: cinco estimadores, seis cenários}
\label{sec:sintese-final}

A Tabela~\ref{tab:rmse-final} reúne o RMSE de $\theta$ para os cinco
estimadores nos seis cenários (malha fechada real, $u=-K\hat X$/\emph{swing-up}
via \texttt{swingup.m}).

\begin{table}[H]
\centering
\small
\begin{tabular}{@{}l r r r r r@{}}
\hline
\textbf{Cenário} & \textbf{Baseline} & \textbf{Complementar} & \textbf{KF} & \textbf{EKF} & \textbf{UKF} \\
\hline
S1 — Estabilização        &  0,58 & 15,58 &  0,24 & 0,25 & 0,24 \\
S2 — \emph{Swing-up}      &  9,75 & 51,62 & 10,51 & 0,24 & 0,23 \\
S3 — Ruído alto            &  9,55 & 16,77 &  1,06 & 1,17 & 1,09 \\
S4 — Baixa taxa/\emph{dropout} &  0,87 & 17,58 &  0,38 & 0,24 & 0,26 \\
S5 — $Q/R$ mal ajustados   &  0,63 & 15,16 &  0,29 & 0,29 & 0,27 \\
S6 — Inicialização ruim    &  1,29 & 15,78 &  0,24 & 0,24 & 0,24 \\
\hline
\end{tabular}
\caption{RMSE de $\theta$ (graus) por estimador e cenário, malha fechada real.}
\label{tab:rmse-final}
\end{table}

Além do já discutido S2:

\begin{itemize}
  \item \textbf{EKF e UKF praticamente empatados} em todos os cenários, já
  que o modelo de medição é linear (Seção~\ref{sec:ukf}): a diferença está só
  em como cada um propaga a predição não-linear, e raramente é decisiva neste
  projeto. A vantagem do UKF aparece no \emph{swing-up} (S2), com pequena
  vantagem sobre o EKF ($0{,}23^{\circ}$ \emph{vs.} $0{,}24^{\circ}$; mais
  nítida em $\dot\theta$, Tabela~\ref{tab:s2-swingup}).

  \item \textbf{Regime quase-linear (S1, S4, S5, S6):} KF, EKF e UKF
  convergem para desempenho parecido, a poucos centésimos de grau um do
  outro --- EKF/UKF não perdem nada por levar um modelo mais caro quando a
  situação não exige.

  \item \textbf{S5 ($Q/R$ mal ajustados) afeta os três "com modelo" igualmente:}
  o NEES médio sobe de forma quase idêntica (KF $19{,}19$, EKF $19{,}66$, UKF
  $18{,}29$, todos acima do limite $\chi^2_4$ de $11{,}14$) --- a
  inconsistência é uma propriedade do ajuste de $Q/R$, não do tipo de
  estimador.

  \item \textbf{Anomalia pontual em S6:} sob inicialização ruim ($P_0=100I$),
  o RMSE de $\dot\theta$ do UKF sobe para $0{,}236$~rad/s (EKF $0{,}026$, KF
  $0{,}045$), embora o RMSE de $\theta$ fique idêntico ($0{,}24^{\circ}$) ---
  provável transiente nos \emph{sigma points} iniciais, espalhados demais
  para a dinâmica RK4 de um passo. Não compromete a conclusão geral; fica
  como candidato a investigação futura.

  \item \textbf{$t_{converge}$ é uma métrica frágil perto do limiar.} Em S3,
  o KF converge em $3{,}26$~s enquanto EKF/UKF aparecem como \texttt{NaN},
  apesar de RMSE de $\theta$ comparável ou melhor nos três. A causa é o RMSE
  de $\dot\theta$ colado na tolerância de convergência ($0{,}10$~rad/s):
  \texttt{compute\_metrics.m} exige os quatro estados dentro da tolerância por
  1~s contínuo, e um único passo acima do limite reinicia a janela ---
  sensível à sorte da realização de ruído perto da fronteira. O mesmo explica
  por que baseline e complementar quase nunca convergem por esse critério.
\end{itemize}

\figplaceholder{results/cpu\_time.png}{Custo computacional médio por
estimador (\texttt{cpu\_time}, escala log), medido dentro de
\texttt{run\_scenario.m} e médio sobre os seis cenários. A ordem reflete
diretamente a complexidade algorítmica de cada método por passo: baseline e
complementar são $O(1)$ (sem álgebra matricial); KF é uma predição/atualização
$4\times4$ fixa; EKF soma o custo de \texttt{jacobian\_f} + \texttt{expm}
($4\times4$) a cada passo; UKF paga por $2n{+}1=9$ integrações RK4 completas
da dinâmica não-linear por passo, além da fatoração de Cholesky.}
{fig:cpu-time}

Em termos absolutos (Figura~\ref{fig:cpu-time}), baseline e complementar
custam ${\sim}1{,}3$--$1{,}5$~ms por execução completa; o KF, ${\sim}6{,}5$~ms
($5\times$); o EKF, ${\sim}42{,}8$~ms ($33\times$); e o UKF, ${\sim}88{,}6$~ms
($68\times$ o baseline, $2{,}1\times$ o EKF --- preço de propagar $9$ pontos
por RK4 em vez de $1$ Jacobiano). Em nenhum cenário deste projeto o custo é
proibitivo, mas a razão EKF/UKF seria o critério decisivo num microcontrolador
de malha rápida, onde o EKF dá quase a mesma precisão por metade do custo por
passo.

\subsection{O que o ruído explica (e o que não explica): revisitando S0}
\label{sec:ideal-rmse}

A Seção~\ref{sec:ideal} definiu S0 (mesma trajetória de S2, mas com
$\sigma_x=\sigma_\theta=0$) como referência sem ruído. Rodando os cinco
estimadores nesse cenário, separamos dois efeitos que o RMSE de S1--S6 mistura:
erro causado pelo ruído de medição \emph{vs.} erro estrutural do método,
presente mesmo com medição perfeita.

\figplaceholder{results/S0\_ideal\_rmse.png}{RMSE de $\theta$ por estimador em
S0 (sem ruído de medição), escala logarítmica. Baseline e complementar
mantêm erro de ordem $10^1$ mesmo sem ruído algum; KF, EKF e UKF colapsam
para $\lesssim 10^{-4\circ}$.}
{fig:ideal-rmse}

O resultado é nítido: com $\sigma=0$, KF, EKF e UKF reduzem o RMSE de $\theta$
a $5{,}6\times10^{-5\circ}$, $1{,}1\times10^{-8\circ}$ e
$1{,}1\times10^{-8\circ}$ respectivamente --- a estimativa colapsa para a
verdade, como esperado de um filtro bem formulado sem incerteza de medição.
Baseline e complementar, porém, não zeram: ficam em $10{,}28^{\circ}$ e
$51{,}75^{\circ}$, praticamente os mesmos valores já vistos em S2 com ruído
nominal ($9{,}75^{\circ}$ e $51{,}62^{\circ}$, Tabela~\ref{tab:rmse-final}).
Ou seja, o erro desses dois métodos durante o \emph{swing-up} não vem do
ruído do sensor --- vem de diferenciar numericamente um sinal que muda rápido
demais para a constante de tempo do passa-baixas (Seção~\ref{sec:baseline-lowpass})
e para a ponderação fixa do complementar (Seção~\ref{sec:complementar}),
regime não-linear em que nenhum dos dois tem um modelo da dinâmica para se
apoiar. Zerar o ruído do sensor não teria ajudado; o problema é estrutural.
É o argumento mais direto deste relatório a favor dos estimadores baseados em
modelo (KF/EKF/UKF, cada um dentro de seu regime de validade) sobre as
alternativas herdadas do curso.

%=================================================================================================
%   SEÇÃO C — FIM
%=================================================================================================

%-----------------------------------------------------------------------------------------------
%   DISCUSSAO / ARQUITETURA DA MALHA
%-----------------------------------------------------------------------------------------------
\section{Arquitetura da malha e discussão}

A Figura~\ref{fig:malha} reúne os módulos num diagrama de blocos da malha
fechada. O laço opera em tempo discreto (passo $\Delta t$): a cada instante o
sensor mede as posições ($z=[x;\theta]$ com ruído/perda), o estimador (KF, EKF
ou UKF) reconstrói o estado completo $\hat{X}$, o controlador de duas fases
gera a força $u$ e a planta não-linear é integrada até o passo seguinte.

\begin{figure}[H]
\centering
\begin{tikzpicture}[
  >=Latex,
  node distance=1.5cm and 1.9cm,
  block/.style={draw, thick, rounded corners, align=center,
                minimum height=1.25cm, minimum width=2.5cm},
  sig/.style={->, thick, line width=0.9pt}
]
  \node[block] (ctrl) {Controlador\\[-2pt]\small swing-up + PD\\[-2pt]\small $\rightarrow$ LQR};
  \node[block, right=of ctrl] (plant) {Planta\\[-2pt]\small cart-pole\\[-2pt]\small (não-linear)};
  \node[block, right=of plant] (sensor) {Sensor\\[-2pt]\small ruído, dropout,\\[-2pt]\small outliers};
  \node[block, below=1.6cm of sensor] (est) {Estimador\\[-2pt]\small KF / EKF / UKF};

  \draw[sig] ($(ctrl.west)+(-1.5,0)$) -- node[above,align=center]{$X^\ast{=}0$}(ctrl.west);
  \draw[sig] (ctrl) -- node[above]{$u$} (plant);
  \draw[sig] (plant) -- node[above,align=center]{$x$\\\small (verdade)} (sensor);
  \draw[sig] (sensor) -- node[right]{$z$} (est);
  \coordinate (fb) at (ctrl.south |- est.west);
  \draw[sig] (est.west) -- node[below]{$\hat{X}$ (estimado)} (fb) -- (ctrl.south);
\end{tikzpicture}
\caption{Malha de controle com estimação de estados. O controlador nunca acessa
o estado verdadeiro: fecha o laço sobre a estimativa $\hat{X}$ (princípio da
separação). É neste ponto que se comparam os cinco estimadores.}
\label{fig:malha}
\end{figure}

Dois aspectos merecem destaque. Primeiro, o princípio da separação: o
controlador foi projetado supondo acesso ao estado, mas opera sobre $\hat{X}$;
a Seção~\ref{sec:sintese-final} quantifica essa degradação --- com baseline ou
complementar, o LQR (estável por projeto, Seção~\ref{sec:matlab}) nunca chega
ao desempenho da malha ideal, enquanto com KF/EKF/UKF no regime válido a malha
real se aproxima da malha com estado verdadeiro. Segundo, o controlador de
duas fases (energy shaping + PD longe da vertical, LQR na janela de captura)
leva os filtros, num único experimento, por um regime fortemente não-linear e
por um quase-linear --- mostrando onde cada técnica de estimação compensa,
mais nítido na divergência do KF e na robustez de EKF/UKF durante o
\emph{swing-up} (Seção~\ref{sec:resultados-naolineares}).

\section*{Conclusão}
\addcontentsline{toc}{section}{Conclusão}

Implementamos e comparamos, dentro de uma malha de controle fechada real,
quatro níveis de estimadores de estado --- diferenciação com passa-baixas,
filtro complementar, KF linear, EKF e UKF --- aplicados ao pêndulo invertido
sobre carro. As conclusões centrais, sustentadas pelos seis cenários (S1--S6)
rodados com o harness completo (\texttt{src/main.m}):

\begin{itemize}
  \item \textbf{O KF linear é ótimo só perto do ponto onde foi linearizado.}
  Empata com EKF/UKF no regime quase-linear (S1/S4/S5/S6); no \emph{swing-up}
  (S2) diverge de forma otimista --- a covariância relatada encolhe
  exatamente quando o erro real explode (NEES $\sim\!2{,}5\times10^{5}$
  contra o limite $\chi^2_4$ de $11{,}14$).
  \item \textbf{EKF e UKF resolvem esse problema}, carregando o modelo
  não-linear correto: RMSE de $\theta$ abaixo de $0{,}25^{\circ}$ e NEES
  dentro da faixa esperada nos seis cenários. A diferença entre os dois é
  pequena aqui; o UKF leva vantagem marginal em regimes fortemente
  não-lineares, ao custo de ${\sim}2\times$ o tempo de CPU do EKF.
  \item \textbf{"Mais modelo" nem sempre é melhor}: o filtro complementar é o
  pior estimador em quase todo cenário, e o baseline sem modelo supera o KF
  divergente no \emph{swing-up}.
  \item \textbf{O princípio da separação se verifica na prática}: o mesmo
  ganho LQR (estável por projeto, Seção~\ref{sec:matlab}) produz malhas de
  qualidade muito diferente dependendo do estimador que alimenta $\hat X$.
  \item Os testes de consistência (NEES/NIS) revelam o que o RMSE esconde: em
  $Q/R$ mal ajustados (S5), o RMSE quase não piora, mas os filtros baseados
  em modelo relatam incerteza menor que a real --- só o teste $\chi^2$ expõe
  isso.
\end{itemize}

Como trabalho futuro, ficam abertos: recalibrar $Q$ por um método rigoroso de
densidade espectral (em vez do ajuste direto em tempo discreto,
Seção~\ref{sec:kf-algoritmo}); adotar a forma de Joseph também na atualização
do UKF (Seção~\ref{sec:ukf}); e investigar a anomalia pontual de $\dot\theta$
do UKF sob inicialização ruim (S6, Seção~\ref{sec:sintese-final}).


\section{Uso de IA}

Usamos nesse projeto o Claude Code, desde a elaboração dos modelos usando matlab até a escrita deste relatório.
Apesar do grande auxílio que essa ferramenta trouxe, não deixamos de aprender e utilizar nosso pensamento crítico quanto às decisões de projeto. A concepção do modelo físico foi feita à mão, as implementações dos filtros foram verificadas e o modelo matemático por trás também, com cada aluno cumprindo uma parte do projeto. 

Obs: esta seção foi escrita por um humano :)
%-----------------------------------------------------------------------------------------------
%   REFERENCIAS
%-----------------------------------------------------------------------------------------------
\begin{thebibliography}{9}

\bibitem{mathworks}
MathWorks. \emph{Derive and Simulate Cart-Pole System} (Symbolic Math Toolbox).
Disponível em:
\url{https://www.mathworks.com/help/symbolic/derive-and-simulate-cart-pole-system.html}.

\bibitem{mit}
R.~Tedrake. \emph{Underactuated Robotics: Algorithms for Walking, Running,
Swimming, Flying, and Manipulation} --- Cap. \emph{The Simple Pendulum}.
Disponível em: \url{https://underactuated.mit.edu/pend.html}.

\bibitem{welch1995}
G.~Welch, G.~Bishop. \emph{An Introduction to the Kalman Filter}. Technical
Report TR~95-041, University of North Carolina at Chapel Hill, 1995
(atualizado em 2006). Disponível em:
\url{https://www.cs.unc.edu/~welch/media/pdf/kalman_intro.pdf}.

\bibitem{simon2006}
D.~Simon. \emph{Optimal State Estimation: Kalman, $H_\infty$, and Nonlinear
Approaches}. Wiley, 2006.

\bibitem{wan2000}
E.~A.~Wan, R.~van der Merwe. \emph{The Unscented Kalman Filter for Nonlinear
Estimation}. Proc. IEEE Adaptive Systems for Signal Processing, Communications,
and Control Symposium (AS-SPCC), pp.~153--158, 2000.

\bibitem{oppenheim}
A.~V.~Oppenheim, R.~W.~Schafer. \emph{Discrete-Time Signal Processing}, 3.ª
ed. Prentice Hall, 2009.

\end{thebibliography}

\end{document}
